\section{Discussion}
\label{sec:discussion}

\subsection{What We Actually Did}

We used real CSI 300 data from Baostock API. The 14.79\% annual return with 15.31\% drawdown is what happened, not what we simulated.

The 60-day momentum factor showed ICIR = 1.49, which is higher than typical developed-market findings. Whether this generalizes beyond our sample is an open question. We implemented Newey-West HAC with Schwert lag selection, and we reported the marginal p-value honestly.

\subsection{Why Momentum Might Work Better in China}

We see three possibilities. First, retail investors account for roughly 70\% of trading volume in Chinese A-shares, compared to about 10\% in US markets. Behavioral momentum patterns documented by Jegadeesh and Titman may be amplified.

Second, information diffusion appears slower in emerging markets, which could extend momentum horizons. The 60-day signal worked better than shorter windows in our tests.

Third, our sample period (2020-2024) lacked the extreme crashes of 2008 or 2015. Momentum may be less effective when markets melt down. We cannot test this with our data.

\subsection{Market Microstructure: Why Monthly Rebalancing Matters}

We chose monthly rebalancing for practical reasons. Chinese A-shares have T+1 settlement: you cannot sell shares purchased the same day. This does not affect us since we hold for 20 trading days. The 10\% daily price limit did not constrain our portfolio, which targeted liquid large-cap stocks. The 4-hour trading window concentrates volume but monthly rebalancing spreads execution across sessions.

These microstructure features informed our design but did not drive it. The monthly frequency balanced signal decay against transaction costs.

\subsection{Transaction Costs Matter}

Round-trip costs averaged 0.15\%: 0.03\% commission, 0.10\% stamp duty on sells, 0.05\% slippage estimate. Monthly rebalancing (12 cycles per year) cost about 1.8\% annually in friction.

Weekly rebalancing would cost roughly 7.2\% annually. Daily would hit 18\%. The alpha would not survive. We tested this: more frequent rebalancing degraded returns.

\subsection{The p-Value Question}

The two-tailed p-value is 0.085. The one-tailed is 0.0425. For trading strategies, the one-tailed test makes more sense: we care whether returns are positive, not whether they differ from zero in either direction.

That said, 60 monthly observations is a small sample. The result could be noise. We report it honestly.

\subsection{Survivorship Bias: The Uncomfortable Truth}

We excluded delisted stocks. This is standard practice but introduces survivorship bias. In US markets, this typically inflates returns by 1-3\% annually. Chinese A-shares have lower delisting rates due to regulatory protections, and CSI 300 constituents are large caps with low delisting probability. Our best guess: the bias inflates our 14.79\% return by roughly 1-2\% annually.

Future work should use CSMAR or Wind data that includes historical constituents. We did not have access for this study.

\subsection{What We Would Do Differently}

Test the full CSI 300 (300 stocks) rather than 94. Validate across markets (US, Europe). Expand beyond 10 factors to the full 25-factor library. Use 10+ years of data rather than 5. The current study is a first step.

\subsection{Replication}

Everything is documented. Data: Baostock API (free, public). Factor formulas: standard academic definitions in the methodology section. Backtest: walk-forward with parameters specified. Statistics: Newey-West HAC with Schwert lags. Code: \url{https://github.com/RomanCohort/Collegium}.
